% ============================================================================
%  Programación y puesta en marcha
% ============================================================================
\section{Programación y puesta en marcha}
\label{sec:programacion}

\subsection{Herramientas}

\begin{table}[H]
  \centering
  \caption{Herramientas de desarrollo y programación.}
  \label{tab:herramientas}
  \small
  \begin{tabularx}{\textwidth}{@{}l l L@{}}
    \toprule
    \textbf{Herramienta} & \textbf{Versión} & \textbf{Observación} \\
    \midrule
    MPLAB X IDE / IPE & v6.20 & Última versión compatible con el PICkit~3: Microchip retiró su soporte desde la v6.25. \\
    Compilador MPLAB XC8 & v3.10 (también v2.46) & Licencia gratuita; ambas versiones compilan el proyecto. \\
    PICkit~3 Programmer & v3.10 & Aplicación independiente, usada con el programador clon. \\
    Proteus & -- & Simulación con L293D en lugar del DRV8833. \\
    Dabble (Android) & -- & En la Play Store figura como \emph{Arduino ESP Bluetooth -- Dabble}, de STEMpedia. \\
    \bottomrule
  \end{tabularx}
\end{table}

\subsection{Proyecto con varias configuraciones}

Un proyecto de MPLAB X une todos sus archivos \texttt{.c} en un solo
\texttt{.hex}, y dos archivos con su propia función \texttt{main()} producen
errores de definición múltiple. Para conservar el código del docente intacto
junto al propio, el proyecto tiene tres configuraciones (\emph{Project
Properties $\rightarrow$ Manage Configurations}); en cada una se excluyen los
archivos que no corresponden (\emph{File Inclusion/Exclusion}):

\begin{table}[H]
  \centering
  \caption{Configuraciones del proyecto en MPLAB X.}
  \label{tab:configs}
  \small
  \begin{tabularx}{\textwidth}{@{}l l L@{}}
    \toprule
    \textbf{Configuración} & \textbf{Archivo} & \textbf{Contenido} \\
    \midrule
    \texttt{try1}     & \texttt{codigo\_docente.c} & Código de la guía (secuencia con rampas y freno). \\
    \texttt{completo} & \texttt{taller2\_robot\_diferencial.c} & Firmware propio (secuencia o control por celular). \\
    \texttt{leds}     & \texttt{prueba\_leds.c} & Prueba de montaje: secuencia en los tres LEDs. \\
    \bottomrule
  \end{tabularx}
\end{table}

Cada configuración genera su propio archivo en
\texttt{dist/<configuración>/production/}. El nombre del \texttt{.hex} es el
mismo en todas; lo que cambia es la carpeta, y el \emph{checksum} que muestra
el panel del proyecto permite confirmar cuál se generó (por ejemplo,
\texttt{0x5EB1} para el código del docente).

\subsection{Grabación con el PICkit 3}

\begin{enumerate}
  \item Conectar el PICkit al conector ICSP con la \textbf{flecha del pin 1}
        hacia la marca de pin 1 de la tarjeta, y los seis pines conectados.
  \item En \emph{PICkit~3 Programmer}: \emph{Tools $\rightarrow$ Check
        Communication}. Si el clon no aparece, cargar el sistema operativo
        \texttt{PK3OSV020005.hex} con \emph{Tools $\rightarrow$ Download
        PICkit Operating System}.
  \item Activar \emph{Programmer $\rightarrow$ Manual Device Select}, elegir
        \emph{Midrange $\rightarrow$ Standard} y el PIC16F877A. Con la
        detección automática la aplicación se bloqueaba con el clon.
  \item Alimentación: con la tarjeta energizada por su batería, \emph{VDD On}
        desactivado; para pruebas sin motores, \emph{VDD On} a
        \SI{5.0}{\volt}.
  \item \emph{File $\rightarrow$ Import Hex}, \emph{Write} y verificar el
        mensaje \emph{Programming Successful}. Desmarcar \emph{/MCLR} para que
        el programa arranque.
\end{enumerate}

\noindent\textbf{Incidente.} En una de las tarjetas el conector se conectó
invertido. Con la inversión, el voltaje de programación (unos
\SI{13}{\volt}) que el PICkit aplica en su pin~1 llegó al pin 6 del conector,
que corresponde a RB3/PGM, una entrada de E/S cuyo máximo es VDD +
\SI{0.3}{\volt}. Desde entonces esa tarjeta identifica el PIC pero falla al
escribir en la dirección 0x0000, mientras que el mismo programador y el mismo
\texttt{.hex} graban sin problema en otra tarjeta. El síntoma es coherente con
un daño en RB3: en un PIC borrado la programación en bajo voltaje viene
activada, y con PGM alterado el chip no acepta la programación normal. La
tarjeta quedó pendiente de cambio del microcontrolador.

\subsection{Conexión del celular}

\begin{enumerate}
  \item Grabar la configuración \texttt{completo} con \texttt{MODO\_BLUETOOTH
        = 1}. Al encender, los tres LEDs parpadean tres veces y el LED del
        HC-05 parpadea rápido.
  \item Vincular el HC-05 en \emph{Ajustes $\rightarrow$ Bluetooth} del
        celular (clave 1234 o 0000; algunos módulos vinculan sin clave). Debe
        quedar \emph{vinculado}, no \emph{conectado}: si el sistema se queda
        con la conexión, el módulo no aparece en la aplicación.
  \item En Dabble, con la ubicación del celular activa y los permisos de
        dispositivos cercanos concedidos, tocar el ícono de conexión, elegir
        HC-05 y abrir el módulo \emph{Gamepad}.
  \item Presionar varios botones hasta que el LED1 empiece a cambiar con cada
        pulsación: la UART encontró la velocidad del módulo.
  \item Antes de apagar el carro, desconectar desde la aplicación para que el
        celular no conserve la conexión anterior.
\end{enumerate}

\noindent\textbf{Limitación con iPhone.} El HC-05 usa Bluetooth Clásico con el
perfil de puerto serie (SPP), que iOS solo permite a accesorios certificados
por Apple. Por eso el módulo no aparece ni en los ajustes del iPhone ni en
aplicaciones como \emph{Bluetooth Serial}. Para usar un iPhone basta con
reemplazar el HC-05 por un módulo BLE (HM-10, o JDY-33 de doble modo): trabajan
a 9600~baudios, entregan los mismos bytes al PIC y Dabble para iOS los admite,
por lo que el firmware no cambia.
